\documentclass{article}

\usepackage{corl_2026} % anonymous submission mode (no [final]/[preprint])
\usepackage{graphicx}
\usepackage{booktabs}
\usepackage{amsmath,amssymb}

\graphicspath{{generated/}}
% Macros (numbers, table snippets) written by voi_rank.analysis.tables (macros.tex) and
% voi_rank.analysis.extra (macros_extra.tex); absent before the analysis runs.
\InputIfFileExists{generated/macros.tex}{}{}
\InputIfFileExists{generated/macros_extra.tex}{}{}
\providecommand{\voiNScenarios}{N}
\providecommand{\voiNAttempts}{M}
\providecommand{\voiRunId}{--}
\providecommand{\voiSeed}{--}
\providecommand{\voiNDraws}{--}
\providecommand{\voiCodeHash}{--}
\providecommand{\voiDataHash}{--}

\title{Does the sim-to-real gap matter for the value of a robot safety evaluation?}

\author{
  Anonymous Author(s)\\
  Affiliation\\
  \texttt{email}
}

\begin{document}
\maketitle

\begin{abstract}
Robot safety evaluations can be run in simulation, on hardware, or both. Simulation is cheap but its result is a noisier signal about real-world behaviour; hardware is faithful but expensive and slow. We ask when the sim-to-real gap changes whether an evaluation is worth running at all. Framing each evaluation as an instrument informing a deployment decision, we estimate its expected value of information (VOI) under elicited priors and instrument reliabilities~\citep{raiffa1961applied}, and sweep the reliability penalty attributable to the sim-to-real gap across \voiNScenarios{} robot safety scenarios. We find that [placeholder: for which classes of decision the gap flips the sim-versus-hardware choice, and where simulation retains most of the value]. We release scenarios, protocols and a reproducible pipeline.
\end{abstract}

\keywords{Sim-to-real, Safety evaluation, Value of information}

\section{Introduction}
[Placeholder: motivation, the sim-versus-hardware evaluation choice as a decision problem, contributions.]

\section{Value-of-information model}
A decision maker chooses an action $a \in \mathcal{A}$ with utility $u(a,\theta)$ under prior $p(\theta)$. An evaluation returns $y \sim p(y \mid \theta)$; a simulated evaluation has a reliability degraded by a gap parameter $\gamma$. The expected value of sample information is
\begin{equation}
  \mathrm{EVSI}(\gamma) = \mathbb{E}_{y}\!\left[\max_a \mathbb{E}_{\theta \mid y,\gamma}\, u(a,\theta)\right] - \max_a \mathbb{E}_{\theta}\, u(a,\theta).
\end{equation}
Details of the fitting and the $\gamma$ sweep are in Appendix~\ref{app:protocol}.

\section{Method}
Scenario catalogue, elicitation protocol, gap parameterisation, Monte Carlo estimation, sensitivity analysis. [Placeholder: filled from \texttt{generated/}.]

\section{Results}
[Placeholder: Figure~\ref{fig:sweep} and Table~\ref{tab:ranking} are generated by the analysis pipeline.]

\begin{figure}[t]
  \centering
  \IfFileExists{generated/fig_by_level.pdf}{\includegraphics[width=\linewidth]{fig_by_level.pdf}}{\fbox{\parbox{0.9\linewidth}{\centering\vspace{2em}generated/fig\_by\_level.pdf not yet produced (voi\_rank.analysis.figures)\vspace{2em}}}}
  \caption{Median EVSI, efficiency and cost of each rung of the sim-to-real ladder (\texttt{attributes.level}), by scenario group; adjacent-rung marginals are in \texttt{generated/level\_uplift.tex} (\texttt{voi\_rank.analysis.extra}).}
  \label{fig:sweep}
\end{figure}

\begin{table}[t]
  \centering
  \caption{Scenarios ranked by VOI under simulation and hardware. Generated by \texttt{voi\_rank.analysis.tables}.}
  \label{tab:ranking}
  \IfFileExists{generated/ranking.tex}{\input{generated/ranking.tex}}{\begin{tabular}{lrr}\toprule Scenario & VOI (sim) & VOI (hw) \\ \midrule (pending) & -- & -- \\ \bottomrule\end{tabular}}
\end{table}

\section{Discussion and limitations}
[Placeholder.]

\label{lastbodypage}
\clearpage

\bibliography{refs}

\appendix
\section{Elicitation protocol}
\label{app:protocol}
[Placeholder.]

\section{Scenario catalogue}
\label{app:catalogue}
\IfFileExists{generated/catalog.tex}{\input{generated/catalog.tex}}{[generated/catalog.tex not yet produced (voi\_rank.analysis.tables).]}

\section{Ablations}
\label{app:ablations}
[Placeholder.]

\section{Additional figures}
\label{app:figures}
[Placeholder.]

\section{Reproduction}
\label{app:repro}
Everything downstream of the elicitation database is deterministic given the stored seed (run \voiRunId: seed \voiSeed, \voiNDraws{} draws, code revision \texttt{\voiCodeHash}, data hash \texttt{\voiDataHash}; a revision suffixed \texttt{-dirty} was run from uncommitted code, the data hash identifies the valid elicitations the run drew from). \voiNAttempts{} elicitation attempts are stored with their prompt hashes and raw responses. See \texttt{README.md} for the run order.

\end{document}
